\documentclass[11pt]{article}
\usepackage[T1]{fontenc}
\usepackage[a4paper,margin=0.92in]{geometry}
\usepackage{amsmath,amssymb,amsthm}
\usepackage{graphicx}
\usepackage{enumitem}
\usepackage[hidelinks]{hyperref}
\usepackage{microtype}

\newtheorem{theorem}{Theorem}
\newtheorem{lemma}[theorem]{Lemma}
\newtheorem{proposition}[theorem]{Proposition}
\theoremstyle{remark}
\newtheorem{remark}[theorem]{Remark}

\newcommand{\C}{\mathbb C}
\newcommand{\R}{\mathbb R}
\newcommand{\N}{\mathbb N}
\newcommand{\M}{\mathcal M}
\newcommand{\NN}{\mathcal{NN}}

\setlength{\parindent}{0pt}
\setlength{\parskip}{0.52em}

\title{Approximate Continuity Restores the Necessary\\
Obstruction for Deep Complex-Valued Networks}
\author{Run \texttt{fa\_banach\_001} --- lane 0}
\date{9 August 2026}

\begin{document}
\maketitle

\begin{center}
\fbox{\parbox{0.91\textwidth}{
\textbf{Status: candidate partial result, likely valid.}
Everywhere approximate continuity is a natural hypothesis strictly weaker
than continuity under which Voigtlaender's necessary condition becomes valid.
Together with his sufficient theorem, it gives a complete characterization
inside this larger class.  No claim of optimality of the hypothesis is made.
}}
\end{center}

\section{Source problem}

Felix Voigtlaender's \emph{The universal approximation theorem for
complex-valued neural networks}, arXiv:2012.03351, studies fixed-depth,
arbitrarily wide complex-valued feedforward networks~\cite{source}.
Let $\NN^d_{\sigma,L}$ denote the resulting scalar-output network class on
$\C^d$ with activation $\sigma:\C\to\C$ and $L$ hidden layers.

The source class $\M$ consists of the locally bounded activations for which the
closure of the set of ordinary discontinuity points has planar Lebesgue
measure zero.  Source Theorem~1.4 proves, for $L\ge2$, that
$\NN^d_{\sigma,L}$ is universal whenever $\sigma\in\M$ satisfies neither
of the following:
\begin{enumerate}[label=(\alph*),leftmargin=2em]
\item $\sigma(z)=p(z,\overline z)$ almost everywhere for some
      $p\in\C[X,Y]$;
\item $\sigma=g$ almost everywhere or $\sigma=\overline g$ almost
      everywhere for an entire function $g:\C\to\C$.
\end{enumerate}
Its converse assumes that $\sigma$ is continuous.  The remark immediately
after the theorem, on arXiv PDF page~4, asks for natural hypotheses weaker
than continuity under which the necessary condition still holds.

\begin{figure}[ht]
\centering
\includegraphics[width=0.83\textwidth]{figures/open_problem_crop.png}
\caption{Source Theorem~1.4 and its future-work question, arXiv PDF page~4.}
\end{figure}

\section{Candidate result}

Write $\lambda$ for planar Lebesgue measure.  A function
$f:\C\to\C$ is \emph{approximately continuous at $z_0$} if, for every
$\varepsilon>0$,
\[
 \lim_{r\downarrow0}
 \frac{\lambda\bigl(\{z\in B_r(z_0):
       |f(z)-f(z_0)|<\varepsilon\}\bigr)}{\lambda(B_r(z_0))}=1.
\]
Thus ordinary neighborhoods are replaced by sets of density one.

\begin{theorem}[complete criterion under approximate continuity]
\label{thm:main}
Let $\sigma\in\M$ be approximately continuous at every point of $\C$.
For every $d\in\N$ and $L\ge2$, the class $\NN^d_{\sigma,L}$ has the
universal approximation property if and only if neither \emph{(a)} nor
\emph{(b)} above holds.

If either obstruction holds, the stronger source conclusion remains true:
there are $F\in C_c(\C^d;\R)$ and $\varepsilon>0$ such that
\[
  \|F-\Phi\|_{L^1(B_1(0))}\ge\varepsilon
  \qquad\text{for every }\Phi\in\NN^d_{\sigma,L}.
\]
\end{theorem}

\subsection*{Proof intuition}

The exceptional values of an almost-everywhere representative matter because
deep networks compose the activation with itself.  Approximate continuity
removes precisely this freedom: a density-one approach to every point sees
the continuous representative, so the value at that point is forced to equal
the representative's value.  Consequently, an activation which is almost
everywhere polynomial, entire, or conjugate-entire is pointwise so.  The
source obstruction proof then applies verbatim.  Its sufficient theorem
already requires only $\sigma\in\M$, so the two directions meet.

\begin{lemma}[density-one rigidity]
\label{lem:rigidity}
Let $f,h:\C\to\C$.  Suppose that $f=h$ almost everywhere, $f$ is
approximately continuous at $z_0$, and $h$ is continuous at $z_0$.
Then $f(z_0)=h(z_0)$.
\end{lemma}

\begin{proof}
Fix $\eta>0$.  Approximate continuity says that
\[
 A_r:=\{z\in B_r(z_0):|f(z)-f(z_0)|<\eta\}
\]
has positive measure for every sufficiently small $r$.  Continuity of $h$
allows $r$ to be chosen so that $|h(z)-h(z_0)|<\eta$ throughout
$B_r(z_0)$.  Since the set on which $f=h$ has full measure, there is a point
$z\in A_r$ with $f(z)=h(z)$.  Hence
\[
 |f(z_0)-h(z_0)|
 \le |f(z_0)-f(z)|+|h(z)-h(z_0)|<2\eta.
\]
Letting $\eta\downarrow0$ proves the claim.
\end{proof}

\begin{proof}[Proof of Theorem~\ref{thm:main}]
Suppose first that (a) holds.  The function
$h(z)=p(z,\overline z)$ is continuous, so Lemma~\ref{lem:rigidity}, applied
at every $z_0\in\C$, gives $\sigma=h$ pointwise.  In case (b), the same
argument applies to $h=g$ or $h=\overline g$.  In every case $\sigma$ is
therefore continuous and satisfies the corresponding pointwise obstruction.
Voigtlaender's Theorem~4.12 now supplies the displayed $L^1$ lower bound,
and in particular rules out universality~\cite{source}.

Conversely, if neither (a) nor (b) holds, source Theorem~4.11 (equivalently,
the sufficient direction of Theorem~1.4) applies directly to
$\sigma\in\M$ and yields universality for every $d\in\N$ and $L\ge2$.
\end{proof}

\section{The hypothesis is strictly weaker than continuity}

The result is not merely continuity under a different name, even after
intersecting with the source class $\M$.

\begin{proposition}
There is a bounded $\sigma\in\M$ which is approximately continuous at every
point of $\C$ but is not continuous.
\end{proposition}

\begin{proof}
Identify the positive real axis with a subset of $\C$.  Put
\[
 a_n=2^{-n},\qquad r_n=2^{-4n-4},\qquad
 D_n=B_{r_n}(a_n) \quad(n\ge1).
\]
The closed disks $\overline D_n$ are pairwise disjoint.  Let
$\psi(w)=\max\{1-|w|,0\}$ and define
\[
 \sigma(z)=\sum_{n=1}^{\infty}
 \psi\!\left(\frac{z-a_n}{r_n}\right).
\]
At most one summand is nonzero at each point.  Away from the origin the disk
family is locally finite, and each bump vanishes on its boundary, so $\sigma$
is continuous on $\C\setminus\{0\}$.  But $\sigma(0)=0$ while
$\sigma(a_n)=1$ and $a_n\to0$; hence $\sigma$ is discontinuous at $0$.

It remains to check approximate continuity at $0$.  Let
$E=\bigcup_nD_n$, and choose $N$ with
$2^{-(N+1)}<R\le2^{-N}$.  No $D_n$ with $n<N$ meets $B_R(0)$, and therefore
\[
 \frac{\lambda(E\cap B_R(0))}{\lambda(B_R(0))}
 \le
 \frac{\sum_{n\ge N}r_n^2}{R^2}
 \le \frac{2^{-6N-6}}{1-2^{-8}}\longrightarrow0.
\]
Since $\sigma=0$ off $E$, every set
$\{z:|\sigma(z)-\sigma(0)|<\varepsilon\}$ has density one at $0$.
Thus $\sigma$ is approximately continuous there, and it is ordinarily
continuous (hence approximately continuous) everywhere else.  Finally,
$0\le\sigma\le1$ and its discontinuity set is exactly $\{0\}$, so
$\sigma\in\M$.
\end{proof}

\begin{remark}
The source's pathological universal activation (Example~4.13) fails the new
hypothesis in the expected place: on the negative real axis its prescribed
point value differs from the approximate limit supplied by the values off the
axis.  Approximate continuity therefore excludes that null-set coding
mechanism without imposing ordinary continuity everywhere.
\end{remark}

\section{Verification, limitations, and novelty check}

The proof has two dependencies: the elementary density-one rigidity lemma
proved above and Voigtlaender's published Theorems~4.11--4.12.  No numerical
or symbolic computation is used.  The main points for human verification are
the all-points quantifier in approximate continuity, the pointwise upgrade in
Lemma~\ref{lem:rigidity}, and the fact that the source's $L^1$ obstruction can
be invoked after that upgrade.

This answers the future-work direction by giving a natural strict weakening
and a complete characterization within it.  It does not characterize all
discontinuous representatives for which non-universality holds, and it does
not assert that approximate continuity is minimal.

A bounded novelty search through 9 August 2026 checked the run indexes; the
exact phrase ``weaker than continuity'' from arXiv:2012.03351; arXiv/web
queries combining complex-valued networks with approximate continuity; and
the nearby later papers arXiv:2303.16813 and arXiv:2305.16910, which cite the
source but work with smooth or continuous activations~\cite{optimal,narrow}.
No direct statement of the extension above was found.  Mathematical
confidence is high; novelty confidence is moderate because the rigidity step
is elementary and could be implicit in the literature.

\begin{thebibliography}{9}
\bibitem{source}
F. Voigtlaender,
\emph{The universal approximation theorem for complex-valued neural networks},
Appl. Comput. Harmon. Anal. \textbf{64} (2023), 33--61;
arXiv:2012.03351. doi:10.1016/j.acha.2022.12.002.

\bibitem{optimal}
P. Geuchen and F. Voigtlaender,
\emph{Optimal approximation using complex-valued neural networks},
Adv. Neural Inf. Process. Syst. \textbf{36} (2023);
arXiv:2303.16813.

\bibitem{narrow}
P. Geuchen, T. Jahn, and H. Matt,
\emph{Universal approximation with complex-valued deep narrow neural networks},
arXiv:2305.16910; later published in Constructive Approximation.
\end{thebibliography}

\end{document}
